\section{Discussion}
\label{sec:discussion}

\paragraph{The bug in the shipped app.} The misplaced loop did not raise an error, and the
request it sent was valid. The only sign of it was a poor or empty first answer, which is easy to
blame on the model. A unit test of the message builder would have caught it, and so would a
benchmark run like ours, which turned ``the first answer is sometimes odd'' into a measured rate.

\paragraph{Where answers are lost.} When the model was given the annotators' evidence, it was
judged correct for \OracleCorrect{} of the questions, compared with \FixedCorrect{} with the app's
retrieval. There is therefore a gap of about fifteen points between what the app retrieves and
what the model could do with the right evidence. However, the retrieval settings that put
evidence in the context more often closed very little of this gap. In the sweep, the best set-up
at the same prompt length retrieved evidence for about sixteen more questions in a hundred, but
neither of the leading set-ups changed the answers measurably. Reranking added evidence on top of
the smaller chunks, and still the answers did not change. Section-aware chunking and retrying
when the model judged its context insufficient did not change the answers either, partly because
the model accepted \RetryMissedKeptRate{} of the contexts that lacked the evidence.

We can think of three reasons for this, and our experiments cannot tell them apart. First, the
oracle context contains only the evidence, while a retrieved context contains the evidence
among other passages, so the model may need the evidence on its own, or in full, and not just
somewhere in the prompt. Second, a 500-character chunk can match a long evidence paragraph while
holding only part of it, so some of the retrieval gain may come from our matching rule
(Section~\ref{sec:res-sweep}). Third, a correct answer does not always need the matched
evidence. \CorrectWhenMiss{} of the app's answers were judged correct without it, either from
nearby text or because the answer said, rightly, that the paper does not answer the question.
For these reasons, retrieval rates were a poor guide to answer quality in our app, and we
recommend choosing retrieval settings on judged answers.

\paragraph{Chunk size and prompt length.} When chunk sizes are compared with the same number of
chunks, two things change at once, the chunk size and the length of the prompt. At the same
prompt length, smaller chunks retrieved evidence more often in our app, and a cross-encoder
reranker added more. On answers, though, the gain was at most a few points and could not be told
apart from no gain. Long contexts also have a cost for the model~\cite{liu2024lost}, which is
why we kept the prompt length fixed rather than increase it.

\paragraph{Several papers.} A comparison question needs the search to cover all loaded papers,
but a question about one paper does not. Letting the user choose which paper to search is a
simple fix, and it is what the app now does. Sending a question automatically to the paper it
names would make this choice unnecessary, but we have not tested it. The app's rule of one chunk
per file gives a comparison question little evidence from every paper except the best-matching
one, and we saw one comparison answer mix up the test sets of the two papers for this reason
(SI, Section~S7). With two chunks per paper, comparison answers were correct for both papers in
\CompareBalancedBothCorrect{} of pairs, compared with \CompareAppBothCorrect{}. This difference
may be due to chance ($p =$ \CompareP{}), but the change costs nothing in prompt length, so the
app now offers it as a compare mode.

\paragraph{How reliable the results are.} Individual verdicts changed between runs much more
than the overall rates did. \RepeatFixedChanged{} of the app's verdicts changed across the three
runs, while its correct rate stayed between \RepeatFixedLow{} and \RepeatFixedHigh{}. A single
run can therefore support a claim about a rate, but not about any one answer. It also means that
a difference of a few points between set-ups needs more than one run, even when a paired test
gives a small $p$-value, as two runs of the same set-up differed by \RepeatWidestPoints{} points
with $p =$ \RepeatWidestP{}. Most of this variation came from the answering model, not the judge.

For the citations, almost all cited pages were pages the model had been given, and we found most
of the cited statements supported by those pages. This means that a user can check an answer
against the pages it cites, and the app now shows those passages under each answer. The check
itself has to be done by a person, because the judge marked as supported every statement that
we marked unsupported.

\paragraph{How far to trust the judge.} The judge was more generous than we were when grading
against the same references, mostly for answers that hedged or gave only part of what was asked.
For someone using the app, our own rate on the sample (\JudgeCheckHumanCorrect{}, or
\JudgeCheckHumanRuleCorrect{} with the rule check, of \JudgeCheckN{}) means that an answer should
be treated as a pointer to the cited pages and not as a final result. For this paper, it means
that the absolute rates are upper bounds. The comparisons between set-ups are paired and graded by
the same judge, so they are less affected if the judge is equally lenient across set-ups.
However, our labels cover only the app's set-up and do not test this. A second labeller on the
same answers would show how much of the gap comes from the judge and how much from our own
grading.

\paragraph{Limitations.} All our results come from one benchmark of NLP papers, with PDFs that
may differ from the versions the questions were written on. We used one answering model in a
4-bit quantization, asked each question without a conversation, and checked the judge against
our own labels on \JudgeCheckN{} answers. The evidence matching is approximate, and its bias
towards small chunks is the reason we tested the leading retrieval set-ups on answers. The
retrieval sweep and the several-papers experiment were exploratory. Finally, our experiments do
not show how the app compares with systems such as PaperQA2~\cite{skarlinski2024paperqa2}, which
use many model calls per question. Running such systems on the same questions is left for
future work.
